% 見本の章（忠実版）．原書どおりに書き起こしたもの（誤植もそのまま）．内容は架空の小冊子のもの．
\chapter{距離}\label{ch:1}
\origpage{1}

\section{定義}\label{sec:1.1}

\begin{definition}\label{thm:1.1.1}
集合 $X$ 上の関数 $d\colon X\times X\to[0,\infty)$ が，任意の $x,y,z\in X$ に対して次の 3 つをみたすとき，$d$ を $X$ 上の\term{距離}という．
\begin{equation}\label{eq:1.1}
\begin{gathered}
d(x,y)=0 \iff x=y,\qquad d(x,y)=d(y,x),\\
d(x,z)\le d(x,y)-d(y,z).
\end{gathered}
\end{equation}
\end{definition}

\begin{proposition}\label{thm:1.1.2}
任意の $x,y,z\in X$ に対して $|d(x,z)-d(y,z)|\le d(x,y)$ が成り立つ．
\end{proposition}
\begin{proof}
明らか．
\end{proof}

\begin{proposition}\label{thm:1.1.3}
$d$ が $X$ 上の距離なら，$d'=\min\{d,1\}$ も $X$ 上の距離である．
\end{proposition}
\begin{proof}
三角不等式だけを示す．$d(x,z)\ge 1$ のときは
\begin{equation}\label{eq:1.2}
d'(x,z)=1\le d'(x,y)+d'(y,z)
\end{equation}
である．実際，右辺の一方が 1 ならよく，どちらも 1 未満なら $d'=d$ だから \beqref{eq:1.1} より右辺は $d(x,z)\ge 1$ 以上になる．
\origpage{2}%
$d(x,z)<1$ のときは $d'(x,z)=d(x,z)\le d(x,y)+d(y,z)$ で，右辺の各項は $d'$ 以下よい．
\end{proof}

また，次が成り立つ:
\begin{equation}\label{eq:1.3}
d'(x,y)\le d(x,y).
\end{equation}

\begin{example}\label{thm:1.1.4}
$\mathbb{R}$ 上で $d(x,y)=|x-y|$ は距離である．
\end{example}

\section{\texorpdfstring{$\varepsilon$}{ε} 近傍}\label{sec:1.2}

\begin{theorem}\label{thm:1.2.1}
$x\in X$ と $\varepsilon>0$ に対し，集合 $B(x,\varepsilon)=\{y\in X\mid d(x,y)<\varepsilon\}$ は開集合である．
\end{theorem}

次の章の定理 \bref{thm:2.1.1} で，命題 \bref{thm:1.1.2} を使う．印のついた文．
